\section{Kiến trúc tổng thể của bộ điều khiển F450}

\begin{figure}[h]
\centering
\resizebox{\textwidth}{!}{
\begin{tikzpicture}[
    block/.style={
        rectangle,
        draw,
        rounded corners,
        align=center,
        minimum width=2.8cm,
        minimum height=1.0cm
    },
    arrow/.style={
        -{Latex},
        thick
    },
    node distance=1.4cm
]

% Main horizontal chain
\node[block] (traj) {Trajectory\\Generator};

\node[block, right=of traj] (pos)
{Position PID\\$x,y \rightarrow \phi_d,\theta_d$};

\node[block, right=of pos] (att)
{Attitude PID\\$\phi,\theta \rightarrow \Delta PWM$};

\node[block, right=of att] (mix)
{Mixer\\$PWM_1,\ldots,PWM_4$};

\node[block, right=of mix] (motor)
{Motor Model\\$PWM \rightarrow T$};

\node[block, right=of motor] (drone)
{Drone Physics\\Isaac Sim};

% Upper control blocks
\node[block, above left=1.2cm and 0.2cm of mix] (yaw)
{Yaw PID\\$\psi \rightarrow \Delta PWM_\psi$};

\node[block, above right=1.2cm and 0.2cm of mix] (alt)
{Altitude PID\\$z \rightarrow PWM_{base}$};

% Main arrows
\draw[arrow] (traj) -- (pos);
\draw[arrow] (pos) -- (att);
\draw[arrow] (att) -- (mix);
\draw[arrow] (mix) -- (motor);
\draw[arrow] (motor) -- (drone);

% Yaw and altitude arrows to mixer
\draw[arrow] (yaw) -- (mix);
\draw[arrow] (alt) -- (mix);

% Feedback line
\draw[arrow]
    (drone.south) |- ++(0,-1.0) -| (pos.south)
    node[pos=0.25, below] {state feedback};

\end{tikzpicture}
}
\caption{Cascade PID control structure}
\label{fig:cascade_pid}
\end{figure}

Phần điều khiển chính của drone F450 được tổ chức theo dạng \textit{cascade
PID controller}. Lớp trung tâm là \texttt{F450AttitudeHold} trong file
\texttt{issac\_attitude\_hold.py}. Lớp này không chứa toàn bộ thuật toán trong
một khối lớn, mà ghép nhiều module nhỏ lại với nhau:

\begin{itemize}
    \item \texttt{PositionHoldPID}: giữ vị trí ngang $x,y$ và tạo
    $roll\_target$, $pitch\_target$.
    \item \texttt{AltitudeHoldPID}: giữ độ cao $z$ và tạo $PWM_{base}$.
    \item \texttt{AttitudeHoldPID}: giữ roll/pitch và tạo
    $\Delta PWM_\phi$, $\Delta PWM_\theta$.
    \item \texttt{YawHoldPID}: giữ yaw và tạo $\Delta PWM_\psi$.
    \item \texttt{QuadXPwmMixer}: trộn PWM base với các PWM hiệu chỉnh - $\Delta PWM$ để tạo
    bốn lệnh motor.
    \item \texttt{MotorModel}: đổi PWM thành lực đẩy, dòng điện, RPM và torque
    phản lực của từng motor.
    \item \texttt{PhysicalPropellerSpinner}: làm quay joint cánh quạt vật lý
    theo RPM mô phỏng.
    \item \texttt{TrackingLogger}: ghi trạng thái và target ra file CSV để
    phân tích sau khi bay.
\end{itemize}

\subsection{Luồng dữ liệu điều khiển}

Trong mỗi bước vật lý của Isaac Sim, controller đọc trạng thái hiện tại của
drone:

\begin{equation}
    x,\ y,\ z,\quad
    v_x,\ v_y,\ v_z,\quad
    \phi,\ \theta,\ \psi,\quad
    p,\ q,\ r
\end{equation}

Trong đó $\phi$ là roll, $\theta$ là pitch, $\psi$ là yaw, còn $p,q,r$ là vận
tốc góc quanh các trục body.

Sau đó luồng điều khiển chính là:

\begin{equation}
    x_d,y_d,x,y,v_x,v_y,\psi
    \longrightarrow
    PositionPID
    \longrightarrow
    \phi_d,\theta_d
\end{equation}

\begin{equation}
    z_d,z,v_z
    \longrightarrow
    AltitudePID
    \longrightarrow
    PWM_{base}
\end{equation}

\begin{equation}
    \phi_d,\theta_d,\phi,\theta,p,q
    \longrightarrow
    AttitudePID
    \longrightarrow
    \Delta PWM_\phi,\Delta PWM_\theta
\end{equation}

\begin{equation}
    \psi_d,\psi,r
    \longrightarrow
    YawPID
    \longrightarrow
    \Delta PWM_\psi
\end{equation}

Các output trên được đưa vào mixer:

\begin{equation}
    PWM_{base},
    \Delta PWM_\phi,
    \Delta PWM_\theta,
    \Delta PWM_\psi
    \longrightarrow
    Mixer
    \longrightarrow
    PWM_1,\ PWM_2,\ PWM_3,\ PWM_4
\end{equation}

Cuối cùng, từng lệnh PWM đi qua motor model để tạo lực đẩy:

\begin{equation}
    PWM_i
    \longrightarrow
    MotorModel_i
    \longrightarrow
    T_i,\ RPM_i,\ Q_i
\end{equation}

Trong đó $T_i$ là lực đẩy của motor thứ $i$, còn $Q_i$ là torque phản lực do
cánh quạt tạo ra.

\subsection{Vì sao dùng cấu trúc cascade}

Drone là hệ động lực học nhiều tầng. Vị trí không thể điều khiển trực tiếp bằng
PWM, vì PWM trước hết tạo lực đẩy, lực đẩy tạo moment và gia tốc, rồi gia tốc
mới làm thay đổi vị trí. Vì vậy project chia bài toán thành nhiều tầng:

\begin{itemize}
    \item Tầng ngoài cùng - outer loop quyết định drone cần nghiêng bao nhiêu để quay về vị
    trí mong muốn.
    \item Tầng độ cao quyết định tổng lực nâng cần tăng hay giảm bao nhiêu.
    \item Tầng attitude làm drone bám theo góc roll/pitch mong muốn.
    \item Tầng yaw giữ hướng quay quanh trục đứng.
    \item Mixer biến các lệnh trừu tượng thành lệnh cụ thể cho bốn motor.
\end{itemize}

Cách chia tầng này giúp mỗi PID có một nhiệm vụ rõ ràng. Position PID không cần
biết chi tiết từng motor; attitude PID không cần biết target vị trí đến từ đâu;
motor model chỉ cần nhận PWM và trả về lực.

\subsection{Chu kỳ điều khiển trong \texttt{on\_physics\_step}}

Hàm \texttt{on\_physics\_step(dt)} là vòng lặp điều khiển được gọi mỗi bước
physics. Các bước chính là:

\begin{enumerate}
    \item Lấy rigid body \texttt{base\_link} từ \texttt{dynamic\_control}.
    \item Đọc pose, vận tốc tuyến tính và vận tốc góc.
    \item Đổi quaternion sang Euler bằng \texttt{quat\_to\_euler}.
    \item Nếu chưa đặt target $x,y$, tự lấy vị trí ban đầu làm target giữ vị
    trí.
    \item Tính target roll/pitch từ position controller.
    \item Tính $PWM_{base}$ từ altitude controller.
    \item Tính yaw correction nếu yaw hold đang bật.
    \item Tính roll/pitch correction từ attitude controller.
    \item Trộn bốn lệnh PWM bằng mixer.
    \item Cập nhật bốn motor model và apply lực đẩy vào Isaac Sim.
    \item Tính torque phản lực yaw, quay cánh quạt vật lý, ghi log và in debug.
\end{enumerate}

Như vậy, Isaac Sim vừa là môi trường vật lý nhận lực từ controller, vừa là cảm
biến trả lại trạng thái cho controller ở bước sau.

\subsection{Các tham số mặc định quan trọng}

\begin{center}
\begin{tabular}{|l|l|l|}
\hline
Nhóm & Tham số & Giá trị mặc định trong code \\
\hline
PWM & $PWM_{min}, PWM_{max}$ & $1100,\ 1940$ \\
\hline
F450 main & $z_d$ & $1.5$ m \\
\hline
F450 main & $PWM_{hover}$ & $1650$  \\
\hline
Position & $K_{p_x},K_{d_x},K_{i_x}$ & $0.55,\ 1.30,\ 0.03$ \\
\hline
Position & $K_{p_y},K_{d_y},K_{i_y}$ & $0.55,\ 1.30,\ 0.03$ \\
\hline
Position & giới hạn gia tốc & $1.20\ \mathrm{m/s^2}$ \\
\hline
Position & giới hạn góc & $6.5^\circ$ \\
\hline
Altitude & $K_{p_z},K_{d_z},K_{i_z}$ & $220,\ 150,\ 15$ \\
\hline
Attitude & $K_{p},K_i,K_d$ roll/pitch & $520,\ 25,\ 130$ \\
\hline
Yaw & $K_{p_\psi},K_{i_\psi},K_{d_\psi}$ & $95,\ 8,\ 28$ \\
\hline
\end{tabular}
\end{center}

Lưu ý: các giá trị này có thể thay đổi trong lúc chạy bằng cách áp đặt giá trị mới vào Script Editor của Isaac Sim. 

\subsection{Các hàm tiện ích chung}

File \texttt{control\_utils.py} chứa các hàm nhỏ nhưng quan trọng:

\begin{equation}
    clamp(x,l,h)=\max(l,\min(x,h))
\end{equation}

Hàm này dùng để giới hạn PWM, góc, gia tốc và tích phân PID trong vùng an toàn.

Hàm \texttt{wrap\_angle} đưa góc về khoảng:

\begin{equation}
    -\pi \leq angle \leq \pi
\end{equation}

Nhờ đó sai số góc luôn lấy theo đường ngắn nhất. Ví dụ, sai số giữa
$179^\circ$ và $-179^\circ$ phải là $2^\circ$, không phải $358^\circ$.

Hàm \texttt{quat\_to\_euler} đổi quaternion của Isaac Sim sang roll, pitch,
yaw. Đây là bước cần thiết vì controller attitude và yaw làm việc trực tiếp với
góc Euler.
